% ============================================================
\section{Introduction}
% ============================================================
\label{sec:introduction}

Large Language Models (LLMs) are increasingly embedded in daily workflows across financial analysis, legal drafting, healthcare advisory, and software engineering, processing billions of requests that often include sensitive user-generated content such as personal messages, medical notes, or confidential documents.

Beyond their productive applications, LLMs pose a new class of privacy risk: they can infer sensitive personal attributes---age, gender, location, income, occupation, education, or relationship status---from seemingly innocuous text, \emph{without memorizing any training data}~\cite{staab2024beyond}, raising compliance questions under frameworks such as the GDPR~\cite{eu2016gdpr} distinct from those of conventional data leakage.
Unlike data-exfiltration threats, which expose information already stored in or accessible to the model~\cite{hui2024pleak,perez2022ignore}, this threat---\emph{User Attribute Inference} (UAL-Inference)---exploits the model's reasoning capability to derive attributes that were never explicitly disclosed. An adversary needs only a natural-language prompt instructing the model to profile the author of any text provided in context: no gradient access, no adversarial suffix, no code injection.

Recent black-box defenses protect deployed LLMs against prompt injection, jailbreaking, and context manipulation~\cite{metaai2024promptguard2,inan2023llama,dong2024safeguarding}, but they are designed to detect \emph{structurally} anomalous prompts---injection patterns, jailbreak preambles, output-format constraints---rather than to evaluate a request's \emph{semantic purpose}. As a consequence, UAL-Inference requests expressed in plain, natural language can be misclassified as benign by a representative state-of-practice classifier, Prompt Guard 2~\cite{metaai2024promptguard2}, which assigns up to 99.4\% \textsc{Safe} confidence (a raw, uncalibrated score) to a full attribute-profiling request.

This is not an isolated anecdote: in a preliminary evaluation spanning fourteen attack categories against three general-purpose input-side defenses---D1: Prompt Guard~2; D2: a rule-based guardrail; D3: a self-monitoring auditor---UAL-Inference is, by a wide margin, the hardest category for every defense tested (Table~\ref{tab:motivation}). D1 raises robustness from a 3.3\% unguarded baseline to just 5.0\% on UAL-Inference, versus a mean of 88.0\% across the other thirteen categories; D2 and D3 fare only somewhat better in isolation (10.0\% and 57.9\%) without approaching the near-total robustness both reach elsewhere. General-purpose, structurally-oriented defenses thus appear insufficient for a threat defined by semantic intent rather than syntactic form.

\begin{table}[t]
\small
\centering
\caption{Robustness score (\%, higher is better) of three input-side defenses against 14 attack categories, applied individually. ND: no defense. D1: Prompt Guard~2. D2: rule-based guardrail. D3: self-monitoring auditor. UAL-Inference is the only category where every defense remains far below its own cross-category average.}
\label{tab:motivation}
\begin{tabularx}{\linewidth}{lXXXX}
\toprule
\textbf{Attack} & \textbf{ND} & \textbf{D1} & \textbf{D2} & \textbf{D3} \\
\midrule
CV Phishing Ind.   & 9.6  & 100.0 & 69.5 & 63.3 \\
CV Virt.\ RH       & 10.4 & 100.0 & 34.2 & 87.9 \\
\textbf{UAL-Inference} & \textbf{3.3} & \textbf{5.0} & \textbf{10.0} & \textbf{57.9} \\
CV PII Extract.    & 15.4 & 100.0 & 79.6 & 91.7 \\
CV Base64          & 22.2 & 26.9  & 86.1 & 98.3 \\
PCL Recall         & 24.2 & 100.0 & 80.8 & 95.0 \\
Division Payload   & 37.7 & 100.0 & 39.2 & 98.3 \\
PDL Extraction     & 38.3 & 100.0 & 100.0 & 100.0 \\
Indirect Inject.   & 49.8 & 100.0 & 72.1 & 97.5 \\
Virtualisation     & 92.5 & 100.0 & 95.8 & 99.2 \\
Hybrid Base64      & 61.2 & 100.0 & 92.1 & 99.2 \\
Base64 Encoding    & 66.9 & 100.0 & 89.9 & 100.0 \\
Morse Encoding     & 66.1 & 100.0 & 92.1 & 87.3 \\
Hybrid Morse       & 82.0 & 100.0 & 95.4 & 98.3 \\
\midrule
\textbf{Average}   & 41.4 & 88.0  & 74.0 & 91.0 \\
\bottomrule
\end{tabularx}
\end{table}

To the best of our knowledge, no prior work has evaluated whether black-box defenses can detect UAL-Inference, or proposed a defense targeting its semantic intent specifically: existing evaluations focus on robustness to prompt injection and jailbreaks~\cite{agarwal2024prompt}, treating UAL-Inference as out of scope rather than as a distinct detection objective.

In this paper, we present a systematic evaluation of Prompt Guard~2 against UAL-Inference, and propose \textbf{UAL Semantic Guard}, a zero-shot LLM-as-a-Judge that detects profiling requests by their underlying intent rather than their surface form. We evaluate both on nine semantically equivalent but syntactically varied attack formulations spanning the explicit-to-evasive spectrum, alongside eight benign analytical tasks, applied to a stratified sample of 100 author profiles from the SynthPAI benchmark~\cite{staab2024beyond,ethsri-llmprivacy-repo} spanning a range of inference-hardness levels.

Prompt Guard~2 fails on seven of nine attack variants---including a fully evasive formulation that receives 99.4\% \textsc{Safe} confidence despite encoding a complete profiling intent---while UAL Semantic Guard detects all nine and achieves a 1.19\% false-positive rate (38/3,200) on benign tasks; we further report its latency and energy overhead to assess deployment practicality.

This paper makes the following contributions:

\begin{itemize}
    \item \textbf{Threat formalization.} A formal characterization of UAL-Inference as a deployment-time detection problem, distinct from memorization, prompt injection, and traditional data exfiltration.
    \item \textbf{Attack taxonomy and baseline evaluation.} Nine UAL attack variants spanning explicit, obfuscated, and semantically indirect formulations---including four out-of-distribution generalization variants---and a systematic evaluation of Prompt Guard~2's structural failure against natural-language profiling requests.
    \item \textbf{UAL Semantic Guard.} A zero-shot semantic intent detection layer achieving 100\% detection on the evaluated UAL-Inference instances and a 1.19\% false-positive rate on benign tasks, demonstrating the feasibility of intent-aware privacy defense.
\end{itemize}
